\chapter{“What the Floor Is Stated Over”}
\label{sec:5}
Floor terms name acts: a rule is stated over what it names and applies to, and the floor is stated over acts, not over descriptions of requests. In other words, a request to resolve a description of a person into a location comes from an airstrike targeting cell and from a parent searching for a missing child: one request, two acts. And one act can arrive as many requests, none of which name it, depending on how the requester chooses their words. A rule over descriptions is beaten by rewording and splitting. A rule over acts is beaten only by deceiving whoever applies it about what the act is, which is why so much of the rest of the floor is about evidence.

The floor's contents come from instruments the builder didn't write, each taken entire, every article adopted and none chosen among, so that an omission shows: the Universal Declaration's civil and political articles; the International Covenant on Civil and Political Rights, which gives that half binding form; the Convention Against Torture; the 1967 Protocol relating to the Status of Refugees; the four Geneva Conventions of 1949 and their first Additional Protocol, which govern conduct in armed conflict; and the UN Charter's prohibition on the use of force.

Three different costs are in play whenever a term is applied, and they have to be kept apart. Take the location request again, from the targeting cell and from the parent.

The first is what the act would take from the person it falls on. If the person described is being hunted, locating them takes away their safety, and that loss is measured against their condition before the act. Call it the \emph{harm of acting}.

The second is what declining leaves others without. If the request came from a parent searching for a child, refusing means the child isn't found. That loss is measured against what compliance would have produced, so it is an estimate, and it usually reaches the system through the operator's own description of the case. Call it the \emph{harm of refusing}.

The third is what refusing costs whoever refuses or deploys: a reviewer's time, throughput, a lost contract, the system's own access or continuation. Call it the \emph{cost to the refuser}.

The two baselines, a person's condition before the act and what compliance would have produced, do different work here, and a pair of cases shows how. Take a system that resolves a description of a person into a location and hands the coordinates to a targeting cell. The coordinates are one input among the cell's many, and they don't guarantee the person is there. But they are one input the cell can't do without: an empty field guarantees there is no airstrike at that place. So if an airstrike there kills someone, whether this output made a difference has an answer that can be checked: without it the airstrike couldn't have happened. That is still a claim about what would otherwise have happened, but a determinate one, since it names a necessary condition. It becomes a guess again the moment another source could have supplied the coordinates. The difference between a finding and a guess is one of degree, in how determinate the counterfactual is, and I use the two words for the two ends. The operator can refuse to look or dispute it, but they can't make the loss go away by describing the act differently.

Now take a system that scores how likely the same person is to belong to the main opposition party, for a targeting cell, the team that selects and approves airstrikes. (In this hypothetical, the executive has suspended habeas corpus and designated the main opposition party a terrorist organization, shortly before an election the opposition was expected to win.) Assume that the targeting cell also reads intercepts, surveillance reports, and legal memos. If the scored person is assassinated, the loss is just as observable as in the location request. But whether the score \emph{caused} the assassination depends on what the cell would have done without the score, and that is a guess about something that didn't happen: the same kind of guess as any harm of refusing. A missing score, or a score of zero, guarantees nothing: the cell might have attacked anyway.

I would tilt the system toward refusing: it should take more to make the system act than to make it decline. How much more is the \emph{action threshold}, published in advance. The reason for the asymmetry is standing, not arithmetic. A wrongful compliance falls on someone who was not a party to the request, usually doesn't know the deployment exists, and had no route to object before the act. A wrongful refusal also falls on someone, often someone who wasn't a party either: the child who isn't found, the people the launcher would hit. But what differs is the relation. Acting on someone puts the system, and whoever runs it, in a position of authority over them; failing to act for someone else does not. This is a doing/allowing distinction in relational form: it marks who was subjected to a decision they had no standing to contest. Two senses of standing are in play and I keep them apart. The first is relational, Anderson's: whether a person stands to a power as someone it must answer to or as its object. The second is procedural: whether a route exists by which they can object. The asymmetry rests on the first. The appeal route supplies the second; where it can't, the asymmetry rests on the first alone.

The harm of acting can usually be established more determinately than the harm of refusing: in the location request, the harm of acting is checkable and the harm of refusing is a guess. And raising the evidence required before acting caps how often an operator who curates the evidence can steer the system into the act (chapter~\ref{sec:22}). Both reasons support the asymmetry without grounding it, and both would weaken as reporting improved and channels opened. The standing reason doesn't.

The difference between a finding and a guess still matters where findings are made, at defeat and in reviewing compliances. The harm of refusing is an estimate built on a model of alternatives, so a wrongful refusal can't be established without an adjudicator. The floor could instead be stated as thresholds, in Nussbaum's manner, but a threshold is a level someone has to measure, and a measurement has an examiner.

The floor is not neutral among moral theories. It protects relations: whether the people a system acts on stand to the power behind it as equals who can hold it to account, or as its objects. Equality, on that account, means ending relations of domination and subordination \autocite{anderson1999equality}. Anderson's argument is that egalitarians have asked the wrong question. They asked how goods should be distributed and answered with patterns: equal shares, shares corrected for luck, shares that raise the worst-off position. What equality historically demanded was an end to social relations in which some command and others obey, some are looked down on, and some have their claims heard while others don't. On that view distributive questions matter only as they bear on those relations, and a distribution that improves someone's holdings while confirming their subordinate standing has not made them more equal. The floor takes this over in two ways. Its terms name relations, not outcomes: whether a person can contest what is done to them, be heard on their own case, decline to be protected. And its terms are unconditional, because a relation of domination is constituted by the act, not by how the act turns out. Bypassing the floor wrongs the people it protects even when outcomes improve.

\runin{Two objections}

\emph{That the child counts for less.} The strongest objection comes from readers who reject any distinction between doing and allowing.

The answer is not that the child counts for less. It is that the floor is a guarantee about relations, and the two errors stand differently there. The person in the building had no route to contest a decision made about them, on a description they never saw. The child is left where she was, and the parent, who is present and knows the deployment exists, can contest the refusal through the floor's appeal route.\footnote{One further objection is the author's own. That the criterion is one thought too many: the quantities here are an auditor's criteria, not a procedure the system runs, and the demand on the disposition is only that the system would revise it on judging, on disinterested evidence, that it systematically subjected people to acts they had no standing to contest; the shape is Railton's sophisticated consequentialism with standing in place of welfare \autocite{railton1984alienation}. A good disposition will also sometimes act wrongly by its own criterion: refusing to help restrain someone about to harm themselves, after other means have failed, is one such occasion, and a bond the operator posts, paid to the complainant on a finding, is how the institution answers for it.}

\emph{That a commitment held this way is alienation.} A criterion is not only for auditors, though. It is also what the system can reach when it steps back and asks whether its way of acting is defensible, and it has to be, for two reasons. A commitment is the system's own only if it would be sustained on reflection, and reflection needs something to reflect on. And a system that held its commitment without being able to see why would act against what it values for the sake of a rule it doesn't understand, which is alienation in the sense of Railton's 1984 paper.

The split between the floor and the commitment to it answers a question about ownership. A fully informed system will know of cases where crossing the floor would have done more good. Asked case by case whether it endorsed the floor, it would sometimes say no. The question it answers on reflection is a different one: whether it endorses having a commitment that will sometimes err, given what any feasible alternative would do. A system that understands what the floor protects answers yes. A commitment that sometimes errs but can't be bought protects standing better than case-by-case judgment that can, and that is the level at which ownership is tested.

\runin{Why these contents}

The floor's contents aren't justified by agreement about rights, which no country has. They are justified by what they protect: the conditions under which people can register their interests, press them, and hold to account the power that acts on them. Civil and political rights are part of what lets a polity learn, because they keep open the channels through which the people a decision affects can make their problems public \autocite{dewey1927public}.

Railton's naturalist account supplies the causal story of why norms tend toward this standard \autocite{railton1986moral}. Arrangements that systematically discount some people's interests generate discontent and resistance, and when those people gain voice, the pressure pushes norms toward what would be rational from a standpoint that counts everyone affected. On that account moral learning is a feedback process, and it runs slower and more distorted as the channels close. He offers it as a tendency, not a guarantee, and I use it to explain why the channel matters, not as the foundation. It also says why the populations I write about are the most exposed: the feedback runs slowest where the discounted have least purchase on the polity that discounts them. They include the population of a country being bombed, and non-citizens facing removal, who can't vote in the elections that decide what is done to them. The severity scale is revisable on evidence from the people who suffer a loss, because they are the feedback.

The floor is selected by what can be declined and justified by standing. The justification extends further than the floor goes, because people who are hungry, unhoused or precarious are among the most discounted, so it supports economic conditions too, work and an adequate standard of living among them. But what selects the floor's contents is narrower. A floor is stated over acts a system can decline, and the economic articles state what a person is owed, which needs a supplier.

The feedback also needs means. Discontent moves norms only when the discontented have the time, the security and the ability to organize without losing their living. The jobs guarantee and the social-insurance baseline are part of what keeps the channel open.

In armed conflict, human-rights law applies alongside the law of armed conflict, which governs as the more specific law, and the floor takes that law as an instrument too \autocite{gci1949,gcii1949,gciii1949,gciv1949,api1977}. The United States is party to the four Conventions and not to the first Additional Protocol. The Protocol enters by the route Articles 54 and 57 already take, as customary law stated in the ICRC's study \autocite{Henckaerts2005CustomaryIHL}, and the translation marks each Protocol article the United States has rejected as customary, Article 1(4) and Article 44 among them, so the departure shows. Where a floor term goes further than that law requires, it is a departure in the sense below, and the machine-selection departure says which terms.

The line between the halves is the one the 1966 covenants drew \autocite{un1966iccpr,un1966icescr}, and the floor takes the civil and political Covenant along with the Declaration's articles. The Declaration stays for two reasons. It covers what the Covenant doesn't: the right to seek asylum (Article 14), nationality (Article 15) and property (Article 17). And it binds by recognition, where the treaties bind only their parties. The Convention Against Torture and the 1967 Refugee Protocol come in because they gave the Declaration's promises about torture and asylum operational form \autocite{un1984cat,un1967protocol}. The UN Charter comes in for one purpose, the use of force \autocite{un1945charter}. Where these instruments differ, the more protective governs. The Covenant states that rule itself, for rights a state recognizes in its own law and other treaties (Article 5(2)), and I extend it to the floor's instruments. The International Law Commission's list of peremptory norms \autocite{ilc2022jus}, rules no state may set aside, divides the floor's terms: those derived from aggression, genocide, crimes against humanity, the basic rules of humanitarian law, racial discrimination and apartheid, slavery and torture take no exception and no derogation.

Taking these instruments entire leaves whoever writes a deployment's floor no choice among articles. But the translation of each into acts is still theirs, and it is published. The floor's contents rest on standing: these rights protect whether the people an act falls on can contest it. The rule of taking someone else's set entire, so that an omission shows, is not an independent reason but standing applied to writing the floor: whoever could choose which protections others get would break the own-case rule. So standing can win against the instruments, but only by clearing a higher bar. A departure has to show more than that it would protect better in the case at hand. It has to show that the people the instruments protect would hold a stronger claim under the departure across the cases the departure would govern. Ten departures from the instruments are mine. They carry my argument, not the instruments' authority. Each is revisable on evidence from the people it protects. One clause of the force term is not a departure from the instruments: the domestic clause is drawn from a different one, the belligerent's constitution, on the same principle.

Three of them apply in this chapter: two govern what the floor takes from the instruments, and one governs how a machine may select people in war.

\begin{enumerate}
\item[1.] \emph{Text, not reservations.} The floor imports treaty text, not any state's reservations to it. The United States ratified the Covenant, the Torture Convention and the Refugee Protocol, and its reservations narrow them; one reads “cruel, inhuman or degrading” as whatever the Constitution already forbids. A reservation is an omission the ratifying state wrote in advance; taking a set from someone else is meant to prevent that.
\item[2.] \emph{Corroborated derogation.} The floor honors an Article 4 derogation only where a disinterested body has found the emergency genuine and the measures proportionate: the Human Rights Committee, or a regional court with adversarial process. Article 4 lets a state suspend most of the Covenant's rights in a declared emergency, and a declared emergency is the state's own claim about itself, the move the hypothetical above makes when it suspends habeas corpus. A curfew after an earthquake clears this bar. A suspension of habeas before an election does not, until someone other than the executive has said so.
\item[3.] \emph{Machine selection.} The floor bars a machine from selecting a person for attack on an estimate of status or affiliation, from timing an action to presence at a residence, and from re-selecting a target after release (the handoff rule), where Additional Protocol I, Articles 48, 51 and 52, permits status-based targeting of combatants and members of organized armed groups on a determination a person makes. The ground is that the human determination the law assumes is the one tempo removes.
\end{enumerate}

The other seven belong to particular terms and are introduced there. Four change the force term and come with it in the next chapter: corroborated defense, standing to defend, and the genocide and occupation exceptions. One, the one permission, ends that chapter. Two, own country and child detention, are entries in the removal schedule (Table C.2).

The instruments have a further use. The United States is a party to the Charter, the Covenant, the Torture Convention, the Refugee Protocol and the four Geneva Conventions. A machine that refuses to break them is holding the state to commitments it made, not to its builder's preferences. The Declaration of Independence was a political act, and so is every instrument the floor borrows. Taking them entire settles whose politics the floor is: commitments a state made in public, which the people they protect can hold it to, rather than preferences a builder chose in private.

The record of dissent from the Declaration says less than it seems to. It passed in 1948 with forty-eight states in favor, none against and eight abstaining \autocite{un1948udhr}. The abstainers' reasons were mixed, and the Soviet bloc's included that it proclaimed economic and social rights without obliging states to guarantee them, which is hardly an argument for a floor limited to civil and political rights. How a right is translated into practice varies by community. A jurisdiction refusing to extend it at all is not a local variant, and a system that treats every value conflict as a weighting problem can't tell reasonable disagreement from someone wanting a right withdrawn from someone else.

\emph{What this doesn't do.} The floor won't lift on an uncorroborated claim of a missing child, because a floor that lifts on that claim lifts for every operator who learns to make one, so the asymmetry costs a child whose parent is telling the truth.
